%%%%%%%%%%  scan idx 305  |  folio 292  %%%%%%%%%%
un
carré cartésien de schémas. On suppose \(Y\) et \(Y'\) quasi-compacts, \(X\) et \(Y'\) tor-indépendants sur \(Y\) (III 1.5), \(f\) propre, et pseudo-cohérent, et la conjecture de finitude (III 2.1) vérifiée pour \(f\) et \(f'\) (par exemple \(Y\) et \(Y'\) noethériens, ou \(f\) projectif). Alors, si \(f\) est parfait (resp. \(h\) de tor-dimension finie), on a, pour \(x\in K^{\bullet}(X)\) (resp. \(K_{\bullet}(X)\)),
\[
 h^* f_*(x)=f'_* g^*(x),
\]
où \(h^*,g^*\) sont définis par (2.7.1) (resp. (2.12.1)), et \(f_*,f'_*\) par (2.12.3) (resp. (2.11.1.1)) (noter qu'en vertu de (III 4.7.1) \(f'\) est parfait si \(f\) l'est, et qu'en vertu de (III 3.5.1) \(g\) est de tor-dimension finie si \(h\) l'est).

\textbf{3.2. Passage à la limite sur \(K^{\bullet}_{\naif}\) et \(K_{\bullet}\)} Soit \((X_i,f_{ij})_{i\in I}\) un système projectif filtrant de schémas quasi-compacts et quasi-séparés, les morphismes de transition étant affines (cf. EGA IV 8.2.2). Posons \(X=\varprojlim X_i\), et notons \(f_i:X\to X_i\) l'homomorphisme canonique.

\textbf{3.2.1.} Les homomorphismes (2.7.2) \(f_i^*:K^{\bullet}(X_i)_{\naif}\longrightarrow K^{\bullet}(X)_{\naif}\) définissent grâce à la formule de transitivité (2.7.3) un homomorphisme
\begin{equation}\tag{3.2.1.1}
 \varinjlim K^{\bullet}(X_i)_{\naif}\longrightarrow K^{\bullet}(X)_{\naif}.
\end{equation}
Je dis que \emph{l'homomorphisme} (3.2.1.1) \emph{est un isomorphisme}. En effet, soit \(E\) un \(\mathcal O_X\)-Module localement libre de type fini. D'après (EGA IV 8.5.2 et 8.5.5), il existe un indice \(i\) et un \(\mathcal O_{X_i}\)-Module localement libre de type fini \(E_i\) tel que l'on ait \(E\simeq f_i^*(E_i)\). En vertu de (EGA IV 8.5.2.5), l'image de \(\cl(E_i)\) dans \(\varinjlim K^{\bullet}(X_i)_{\naif}\) ne dépend pas du choix de \(E_i\), notons-la \(u(E)\). D'après (EGA IV 8.5.2 et 8.5.6) l'application
\(u:\Ob\Loclib(X)\longrightarrow \varinjlim K^{\bullet}(X_i)_{\naif}\) ainsi définie est additive (1.5)~;
